\begin{lemma}[Characteristic-independent critical polar coordinate]\label{mod:lamprecht-criticalpolarcoordinate}
Let $F$ be a nonarchimedean local field, let $\psi$ be a nontrivial
additive character, let $m(\chi)=2d+1>1$, and let $\Gamma$ be admissible,
so
\[
 \operatorname{ord}_F(\Gamma)=2d+1+n(\psi).
\]
Choose $\delta\in F^\times$ with $\operatorname{ord}_F(\delta)=d$, and
fix a nontrivial additive character $\psi_0$ of $k_F$.  At the minimal
stationary depth $d+1$, the stationary numerator is the quotient class
\[
 \operatorname{St}_{\Gamma,d+1}(\chi)
 \in \mathcal O_F/\mathfrak p_F^d.
\]
Multiplication of a residue variable by $\delta^2$ descends to the other
side of the stationary pairing.  Pairing these two quotient classes defines
an intrinsic nontrivial additive character
\[
 \mathfrak a_{\chi,\delta}(z)
  =\psi\!\left(\frac{\beta\delta^2\widetilde z}{\Gamma}\right),
 \qquad z\in k_F,
\]
where the displayed value holds for every integral lift $\widetilde z$ and
every permitted stationary representative $\beta$.  Thus there is a unique
coefficient
\[
 A=\kappa_\chi(\delta)\in k_F^\times,
 \qquad \mathfrak a_{\chi,\delta}(z)=\psi_0(Az).
\]
The Lean definition of $A$ has no stationary-representative argument.

For an explicit conductor-one denominator $\gamma_0$ defining $\psi_0$,
Lean also proves, for every permitted $\beta$, the equality in the residue
field
\[
 A=
 \overline{\frac{\beta\delta^2\gamma_0}{\Gamma}}\in k_F.
\]
The quotient is integral of valuation zero.  This is a residue-class
identity only: no equality between arbitrary representatives in $F$ is
asserted.  In particular, the reductions obtained from any two permitted
stationary representatives agree.

If $u\in\mathcal O_F^\times$, $\lambda=\bar u$, and
$\delta'=u\delta$, then
\[
 \mathfrak a_{\chi,\delta'}
   =\mathfrak a_{\chi,\delta}.\operatorname{mulShift}(\lambda^2),
 \qquad
 \kappa_\chi(\delta')=\lambda^2\kappa_\chi(\delta).
\]
The formula depends only on $\lambda$, hence is independent of the chosen
integral unit lift.  If a residual automorphism $\sigma$ fixes $\psi_0$,
precomposition by $\sigma$ sends $A$ to $\sigma^{-1}(A)$, while
precomposition by $\sigma^{-1}$ sends it to $\sigma(A)$.  These are the
manuscript's Frobenius and inverse-Frobenius directions.

For a supplied stationary representative $\beta$, define the complete
critical function
\[
 \Phi_\beta(x)=
 \psi\!\left(\frac{\beta\delta\widetilde x}{\Gamma}\right)
 \chi(1+\delta\widetilde x)^{-1}.
\]
It is independent of every integral lift and satisfies the positive-polar
identity
\[
 \Phi_\beta(x+y)
 =\Phi_\beta(x)\Phi_\beta(y)\psi_0(Axy).
\]
In particular the sign is $+Axy$.

If $\beta'=\beta+c_0$ is another permitted representative, Lean bundles
$c_0=\beta'-\beta$ at exactly the depth $\mathfrak p_F^d$ and defines
$\Delta B\in k_F$ by
\[
 \psi\!\left(\frac{c_0\delta\widetilde z}{\Gamma}\right)
   =\psi_0(\Delta Bz).
\]
It also defines, without choosing a new representative,
\[
 a=\overline{\frac{\beta'-\beta}{\beta\delta}}\in k_F
 \quad\text{and proves}\quad
 \Delta B=Aa.
\]
These are distinct quantities.  The exact representative-translation
formulas are
\[
 \Phi_{\beta'}(x)=\Phi_\beta(x)\psi_0(+\Delta Bx)
 =\Phi_\beta(x)\mathfrak a_{\chi,\delta}(a x).
\]
Thus the polar coefficient is unchanged and only the affine character is
translated.

Finally, Lean packages any nowhere-zero function with this positive polar
law as \texttt{CriticalPolarFunction}.  Its scalar and automorphism
precompositions transport the polar coefficient as
\[
 A\longmapsto\lambda^2A,
 \qquad A\longmapsto\sigma^{-1}A,
\]
respectively.  In odd characteristic its quotient by the standard quadratic
refinement is an additive character, giving a unique $B\in k_F$ such that
\[
 \Phi_\beta(x)=\psi_0(Ax^2/2+Bx).
\]
For this same package,
\[
 B\longmapsto B+C\quad\text{under affine translation by }C,
 \qquad
 B\longmapsto\lambda B,
 \qquad
 B\longmapsto\sigma^{-1}B.
\]
For the local critical functions the representative formula specializes to
$B'=B+\Delta B=B+Aa$.  Their complete finite sum, with no elementary
translation discarded, has exactly the certified phase
\[
 \operatorname{phase}\!\left(\sum_{x\in k_F}\Phi_\beta(x)\right)
 =\operatorname{quadraticPhase}(\psi_0,A,B).
\]
\LeanFile{LanglandsFirstMainLemma/Lamprecht/CriticalPolarCoordinate.lean}
\lean{LanglandsFirstMainLemma.finiteAddCharCoefficient}
\lean{LanglandsFirstMainLemma.finiteAddCharCoefficient_mulShift}
\lean{LanglandsFirstMainLemma.finiteAddCharCoefficient_precompose}
\lean{LanglandsFirstMainLemma.finiteAddCharCoefficient_precompose_symm}
\lean{LanglandsFirstMainLemma.CriticalPolarFunction}
\lean{LanglandsFirstMainLemma.CriticalPolarFunction.translate}
\lean{LanglandsFirstMainLemma.CriticalPolarFunction.scale}
\lean{LanglandsFirstMainLemma.CriticalPolarFunction.precompose}
\lean{LanglandsFirstMainLemma.CriticalPolarFunction.affineCoefficient}
\lean{LanglandsFirstMainLemma.CriticalPolarFunction.eq_quadratic}
\lean{LanglandsFirstMainLemma.CriticalPolarFunction.affineCoefficient_translate}
\lean{LanglandsFirstMainLemma.CriticalPolarFunction.affineCoefficient_scale}
\lean{LanglandsFirstMainLemma.CriticalPolarFunction.affineCoefficient_precompose}
\lean{LanglandsFirstMainLemma.CriticalPolarFunction.affineCoefficient_precompose_symm}
\lean{LanglandsFirstMainLemma.CriticalPolarFunction.phase_sum_eq_quadraticPhase}
\lean{LanglandsFirstMainLemma.criticalPolarVariable}
\lean{LanglandsFirstMainLemma.criticalPolarVariable_integral_lift}
\lean{LanglandsFirstMainLemma.criticalPolarAddChar}
\lean{LanglandsFirstMainLemma.criticalPolarAddChar_integral_lift}
\lean{LanglandsFirstMainLemma.criticalPolarAddChar_ne_one}
\lean{LanglandsFirstMainLemma.criticalPolarCoefficient}
\lean{LanglandsFirstMainLemma.criticalPolarCoefficient_integral_lift}
\lean{LanglandsFirstMainLemma.criticalPolarCoefficient_eq_residue}
\lean{LanglandsFirstMainLemma.criticalPolarResidue_representative_independent}
\lean{LanglandsFirstMainLemma.criticalPolarAddChar_scaleCoordinate}
\lean{LanglandsFirstMainLemma.criticalPolarCoefficient_scaleCoordinate}
\lean{LanglandsFirstMainLemma.criticalPolarCoefficient_scaleCoordinate_lift_independent}
\lean{LanglandsFirstMainLemma.criticalPolarCoefficient_frobenius}
\lean{LanglandsFirstMainLemma.criticalPolarCoefficient_frobenius_symm}
\lean{LanglandsFirstMainLemma.criticalAffineTranslationCoefficient}
\lean{LanglandsFirstMainLemma.criticalStationaryRepresentativeDifference}
\lean{LanglandsFirstMainLemma.criticalRepresentativeTranslationParameter}
\lean{LanglandsFirstMainLemma.criticalAffineTranslationAddChar_eq_polar_mulShift}
\lean{LanglandsFirstMainLemma.criticalRepresentativeAffineCoefficient_eq_polar_mul_parameter}
\lean{LanglandsFirstMainLemma.criticalPolarFunction}
\lean{LanglandsFirstMainLemma.criticalPolarFunction_integral_lift}
\lean{LanglandsFirstMainLemma.criticalPolarFunction_positivePolar}
\lean{LanglandsFirstMainLemma.criticalPolarFunction_changeRepresentative}
\lean{LanglandsFirstMainLemma.criticalPolarFunction_changeRepresentative_intrinsic}
\lean{LanglandsFirstMainLemma.criticalPolarData}
\lean{LanglandsFirstMainLemma.criticalAffineCoefficient}
\lean{LanglandsFirstMainLemma.criticalPolarFunction_eq_quadratic}
\lean{LanglandsFirstMainLemma.criticalPolarFunction_phase_eq_quadraticPhase}
\lean{LanglandsFirstMainLemma.criticalAffineCoefficient_changeRepresentative}
\lean{LanglandsFirstMainLemma.criticalAffineCoefficient_changeRepresentative_eq_add_polar_mul_parameter}
\leanok
\SourceText{Lemma \texttt{lem:critical-polar-coordinate}.}
\uses{mod:lamprecht-stationaryclasscalculus,mod:finitefield-quadraticphase,mod:delta-residueformula}
\FormalizationContract The stationary numerator remains a lattice-quotient
class until an arbitrary representative is explicitly quantified.  Every
lift formula is proved for all permitted integral lifts.  Coordinate and
Frobenius directions, the positive polar sign, the distinction between
$\Delta B$ and $a$, and the odd $A/2$ convention are literal.  No canonical
stationary representative, local-field equality inferred from residue
equality, later parameter theorem, or later case theorem is used.
\end{lemma}
